% !TeX root = ../../../main.tex

Des Weiteren werden dort Funktionen ausgeführt, um die Datenbank vorzubereiten.
Denn für jede Fehlerklasse wird die Datenbank neu aufgesetzt, damit alle zu gleichen Bedingungen starten.
Dafür wird die Datenbank vor Beginn der Experimente einem \emph{WARM-UP} unterzogen.
Dabei werden relevante Daten in die Caches der Datenbank geladen, wodurch nachfolgende Zugriffe schneller
verarbeitet werden können.
Damit die Ergebnisse durch vorherige Experimente nicht verfälscht werden, wird nach jedem Durchlauf eines
Experiments ein \emph{soft reset} der Datenbank durchgeführt.
Während dieser Zeit soll sich die Datenbank von Konflikten befreien, Locks lösen und wieder in den Idle Zustand gehen.
Dafür werden zunächst alle noch vorbereiteten Transaktionen der Datenbank mittels \texttt{ROLLBACK PREPARED}
zurückgesetzt.
Anschließend wird durch \texttt{sleep(2)} eine Wartezeit von zwei Sekunden eingehalten, damit sich die Datenbank
vollständig stabilisieren kann.
Nachdem die Datenbank vorbereitet ist, werden die Experimente ausgeführt.
Dafür wurden vier Experimente implementiert.
\texttt{experiment\_0} dient dabei als Baseline-Experiment und bildet den Ausgangspunkt für die weiteren Experimente.
Es beschreibt den Zustand des Systems ohne den Einsatz eines Retry-Mechanismus und ermöglicht somit einen Vergleich
mit den nachfolgenden Experimenten.
Dies wird für die fünf Concurrency-Werte getestet die in Ausschnitt~\ref{lst:params} gezeigt sind.
Das nächste Experiment, \texttt{experiment\_1}, soll den Einfluss von Retry ohne Delay testen.
Dadurch soll gezeigt werden, ob Retry ohne Delay sinnvoll ist und positive Ergebnisse liefert.
Hier ist insbesondere der \texttt{RETRY\_COUNT} Parameter, im folgenden nur noch Count genannt, wichtig, da dieser
anzeigt, wie sich das Verhalten ändert, wenn mehr Retries erlaubt sind.
Der Count von null wird in diesem Fall nicht beachtet, da dieser redundant ist.
In \texttt{experiment\_2} soll der Einfluss von Delay getestet werden.
Dafür wird der Count auf einen festen Wert gesetzt, hier Index 2 also Count = 3, damit dieser die Ergebnisse nicht
verfälscht und die \emph{immediate Retry} Strategie gewählt, weil die den Einfluss von Delay am besten zeigt, denn
der ist weder zufällig noch anderweitig veränderbar während der Strategie.
Der Delay ist nur in seinem Initialwert variabel, wobei auch der Delay von null nicht beachtet wird.
Als Letztes führt \texttt{experiment\_3} einen Vergleich der Strategien durch.
Es werden die beiden Strategien \emph{Retry with Delay} und \emph{Retry with backoff + jitter Delay}, mit einem festen
initial Delay und variabler Anzahl an Retries getestet.
Der initiale Delay ist auf 0,01 Sekunden festgelegt.
Dadurch lässt sich ein Vergleich zwischen allen Strategien ziehen, da bei allen die Anzahl der Retries variiert
wird.
Für die beiden Strategien, die einen Delay verwenden, wird mit 0,01 Sekunden ein bewusst geringer Delay gewählt,
sodass dieser keinen wesentlichen Einfluss auf die Vergleichbarkeit der Ergebnisse hat.
Da alle Experimente mit allen Concurrency-Werten durchgeführt werden, lässt sich in der Analyse herausfinden, ab
welcher Anzahl an gleichzeitigen Anfragen eine Retry Strategie
kippt.

Die Experimente rufen die Funktion \texttt{\_start\_experiment()} auf und übergeben die Indizes der Parameter.
Die aufgerufene Funktion liest die Werte der Parameter über die Funktion \texttt{\_get\_params()} anhand
ihrer Indizes aus.
Zusätzlich wird anhand der Parameter der \texttt{run\_key} erzeugt, der für die Sortierung in der Analyse benötigt
wird.
Die Parameter und der Schlüssel werden an \texttt{\_start\_experiment()} zurückgegeben.
Nachdem die Parameter und der Schlüssel gesetzt wurden, wird der erste Logeintrag für das Experiment erzeugt.
Die \emph{run\_meta} beschreibt genauer, worum es sich in dem Experiment handelt und wann es gestartet wurde.
Daraufhin wird \texttt{\_run\_batch()} aufgerufen, welche den Log-path anhand des caller\_name setzt und mit Threads
die Anfragen an die Datenbank startet.
Für die parallele Ausführung der Clients wird ein \texttt{ThreadPoolExecutor} mit der durch den Concurrency-Wert
angegebenen Anzahl an Threads verwendet.
Für jede Aufgabe wird ein \texttt{Future}-Objekt erzeugt und in einer Liste gespeichert.
Nachdem die Aufgaben abgeschlossen sind, werden die zugehörigen Ergebnisse über \texttt{as\_completed} nacheinander
verarbeitet.
Dadurch erfolgt das Schreiben in die Logdatei über den aufrufenden Thread, wodurch gleichzeitige Zugriffe auf die
Logdatei durch die Worker-Threads vermieden werden.
Für jeden ausgeführten Vorgang werden die Parameter \texttt{type}, \texttt{run\_id}, \texttt{iteration},
\texttt{retries}, \texttt{attempts\_ms}, \texttt{total\_ms}, \texttt{status}, \texttt{pg\_code}, \texttt{tx\_start}
und \texttt{tx\_finish} erfasst und im JSONL-Format in der Logdatei gespeichert.
Die Iteration wird nur wegen der Information mit geloggt, hat aber keinen weiteren Nutzen.

Im Client wird anhand der übergebenen Strategie entschieden, welche Strategie aufgerufen wird.
Das geschieht durch \texttt{match ... : case ... :} und in den einzelnen Cases wird der Rückgabewert der Strategie
gespeichert.
Die Datenbankverbindung wird ebenfalls in dieser Funktion aufgebaut, da jeder Client eine eigene Verbindung benötigt.
Die Ergebnisse werden dann an die zuvor erwähnte Funktion zurückgegeben.